\chapter{The Company Is Other People}
\chaptermark{The Company Is Other People}

An owner of five companies describes his central constraint plainly. If he
steps away for an hour or two, something can change. Partners have made that
burden manageable. With the right team, he says, the work can be done.

That answer is correct as far as it goes. Every new carrier of work also
acquires exposure: a worker waits for payroll, a supplier waits for payment, a
partner signs a guarantee, and a manager answers for a decision. A team does
not create several copies of a founder. Each person brings judgment, interests,
limits, obligations, and a life that does not belong to the company. Scale
depends on coordinating those contributions without hiding power or dissolving
responsibility.

This is the human architecture of a company. It determines who may decide, who
must be heard, who sees the numbers, who carries the danger, who receives the
upside, how disagreement ends, and what happens when trust fails. A weak
architecture can make talented people look incapable. A sound one lets
ordinary people combine their strengths without requiring anyone to become
superhuman.

\section{Hire for a missing contribution}

Vic Keller, who reports selling nine companies, describes an attractive
company through two linked features: a strong value proposition and unusually
capable people who can share in the wealth they create. He says he has spent
his career hiring people smarter than himself. A Dallas heating-and-air-
conditioning owner condenses the same idea into the advice he received from a
former corporate executive: hire your weakness.

That advice is useful only after the weakness is made specific. ``We need great
people'' is not a hiring plan. Observe where a week actually loses judgment:
founder interruptions, queues, rework, waiting customers, unsafe improvisation,
and decisions nobody can close. Then name the outcome the organization cannot
reliably produce. Is it accurate forecasting, safe field work, enterprise
selling, product judgment, production control, customer repair, or the
management of twenty people rather than five? Name the evidence that would show
a person can carry it.

A young e-commerce founder describes the transition from doing roughly one
million dollars alone to building a team for a larger company. His language is
to duplicate himself three or four times and place smarter specialists around
him. The practical insight is specialization; the misleading part is
duplication. A finance leader, operator, or salesperson should not become a
copy of the founder. The company needs the independent judgment the founder
lacks, including the willingness to say that the founder's idea will not work.

Write a bounded responsibility before a job description. It has eight entries:

\begin{itemize}[leftmargin=0pt,label={},labelsep=0pt,itemsep=0.16em,
  parsep=0pt,topsep=0.35em]
\item \textbf{Outcome.} Name the result that must become more reliable, for
  whom, and how it serves the whole promise.
\item \textbf{Authority.} State which choices, spending, and commitments the
  person may close and which remain reserved.
\item \textbf{Resources.} Supply the time, people, tools, access, training,
  and cooperation the outcome honestly requires.
\item \textbf{Limits.} Preserve safety, law, quality, privacy, labor,
  customer, and solvency boundaries that urgency cannot waive.
\item \textbf{Information.} Give the context, evidence, history,
  dependencies, uncertainty, and direct customer signals needed for judgment.
\item \textbf{Support.} Name the backup, specialist help, and escalation path
  available when the conditions exceed the role.
\item \textbf{Review.} Set the evidence and cadence for examining outcomes,
  complaints, assumptions, workload, and the fit of the role itself.
\item \textbf{Remedy.} State who may stop harmful work and how a failed
  decision repairs the customer, protects people, and improves the system.
\end{itemize}

Suppose a growing service company needs a field lead because the founder still
resolves every exception. \emph{Own the route} is too vague. The lead might be
accountable for accepted work delivered safely, with authority to reschedule
crews and pause a hazardous visit within a stated budget. The company must
supply staffing, tools, contract terms, complaint history, and a reachable
specialist. The lead cannot invent a customer claim or waive safety. A weekly
review examines acceptance, rework, waiting time, worker load, and unresolved
exceptions. If harm occurs, the remit includes immediate protection, customer
repair, preserved facts, and correction of the route. Task, authority, truth,
support, and consequence now travel together.

One medical-device distributor offers the familiar rule to hire slowly and
fire quickly. Speed is a poor substitute for care in either direction. A
company should investigate a consequential hire before making promises, then
state expectations, coach against observable gaps, document material issues,
provide a fair chance to respond, and follow employment law. Immediate removal
is appropriate when safety, violence, fraud, severe harassment, or another
serious duty requires it. Otherwise, the purpose is neither delay nor drama;
it is a timely, evidence-based decision that respects the person and protects
the work.

\section{Test contribution before dividing ownership}

Employment can be ended under defined conditions. Partnership entangles
capital, authority, information, reputation, and exit. Friendship and enthusiasm
do not answer what happens when one person contributes less, needs cash, wants
to sell, becomes ill, changes direction, or believes the other has violated
the agreement.

An Austin real-estate founder describes a more useful path to trust. Before
forming a broad partnership, he gave a future partner small operating tests:
knock on ten doors, make a set of calls, complete concrete work. He reports
that the other man repeatedly did much more than requested. Over two or three
years they accumulated evidence, including several properties owned together,
before dividing a larger company by complementary roles. One became responsible
for operations, profitability, and personnel; the other concentrated on
relationships and new opportunities.

The lesson is not that excess effort proves permanent character. It is that
reversible work reveals more than a partnership pitch. Test how the parties
handle an ordinary deadline, an unpleasant task, a lost customer, a small
amount of money, a correction, and credit for another person's work. Increase
commitment only after the evidence improves.

Glenn Boyd recalls that he and his co-founder argued constantly in their early
company, sometimes over matters as small as a shade of packaging. They learned
which disagreements mattered, developed confidence in each other's judgment,
and eventually made independent decisions. Constructive conflict was part of
their trust formation. The same conflict without decision rules could have
consumed the business.

Another founder reports shutting a company because too many partners wanted
different things. His rule of thumb about the ideal number is personal
experience, not a general limit. His more durable test is that each partner
should make an essential contribution without holding the company captive.
That requires more than complementary personalities. Before issuing equity,
extend the six-field ownership claim from Chapter 10. First record the purpose,
each continuing contribution, and the reversible work that supports a larger
commitment. Then record the economics: equity, vesting or repurchase,
compensation, expenses, distributions, and future capital obligations. Next
set independent and joint decision domains, spending limits, and access to
accounts, contracts, forecasts, and material risks. Finally settle conflicts
of interest, related-party dealings, confidentiality, intellectual property,
outside work, deadlock, incapacity, death, misconduct, departure, valuation,
transfer, and sale. Governing documents, insurance, and qualified advice make
the agreement real.

These terms support trust by answering questions while the parties still
agree. They also prevent a broken relationship from turning uncertainty into
private leverage.

\section{Let many people think and one person close the decision}

The chief executive of a neurotechnology company contrasts vertical managers
who issue orders with horizontal managers who sit together and work through a
problem. He reserves rank for a critical moment in which a final decision must
be made. The distinction is valuable if its limits are explicit. Collaboration
without closure becomes delay; emergency authority without a narrow trigger
becomes ordinary authoritarianism.

A restaurant executive makes the operating point more precisely. Several
people may share responsibility for the work, but one person must remain
accountable for ensuring that it is completed. Accountability here does not
mean absorbing blame for choices made elsewhere. It means holding the final
authority, information, and duty for a named outcome.

For every consequential recurring decision, record five lines:

\begin{description}[leftmargin=2.75em,style=nextline]
\item[Propose] Who frames the choice and supplies the evidence?
\item[Consult] Whose knowledge, rights, or exposure requires participation
  before a decision?
\item[Decide] Which one role closes the decision within the stated boundary?
\item[Execute] Who carries each commitment, with what resources and deadline?
\item[Stop] Who may pause the work for safety, law, solvency, or serious customer
  harm, and how is the stop reviewed?
\end{description}

The record should be proportionate. Choosing office food does not need a
committee. Entering a new country, changing a medical claim, committing scarce
cash, terminating a senior employee, or placing a worker in danger does.
Information rights should follow exposure: a person cannot responsibly carry
an outcome while material facts are withheld.

Critical authority also needs an expiry. Name the event that activates it, the
decisions it permits, the people who must be informed, the evidence retained,
and the condition that returns the company to ordinary governance. Otherwise
every impatient executive can redefine inconvenience as crisis.

A decision map must also follow the people who carry its consequence. Trace
the customer promise, worker exposure, supplier and creditor dependence,
owner capital, regulatory duty, community effect, and anyone placed at
material risk without being a customer. Durable shareholder return is a
legitimate purpose: people who supply risk capital deserve the information,
rights, and possible return they were promised. It does not create automatic
priority over every other claim. A safety right, informed consent, a debt
covenant, a supplier invoice, and a residual equity claim have different
grounds and cannot be collapsed into one vote.

This tracing does not turn leadership into consensus theater. Office food may
still be chosen quickly; a person who lacks relevant knowledge, right, or
exposure need not join every meeting. Consultation widens with consequence,
irreversibility, and uncertainty. One named role closes the decision, while no
role gains the power to make an unconsented loss disappear from the record.

\section{Culture is what the company protects under pressure}

A Texas executive traces customer care back through employee experience. When
people have meaningful work, a safe environment, and reason to believe the
company values them, he argues, they remain close enough to the customer to
generate better ideas. When their attention is consumed by finding the next
paycheck, the company loses that knowledge.

This is more concrete than posters about values. Culture is the pattern of
decisions a new employee can observe: who receives information, whose warning
slows a launch, who gets credit, what a manager is allowed to conceal, whether
an inconvenient customer is repaired, and whether revenue excuses conduct
that would otherwise be unacceptable.

A hospitality operator reports firing his highest-producing salesperson after
the person lied to a customer and damaged the organization. He says the company
then grew by twenty percent. That sequence does not prove the firing caused
the growth, and it should not become a license for impulsive termination. It
does show the governance problem created when one person's revenue purchases
immunity. A star who can violate the promise teaches everyone else that stated
values are optional prices paid by weaker employees.

Honest dissent belongs to the same system. Boyd looks for people who are
intelligent, honest, and able to receive constructive feedback, explicitly
including himself. A safe team is not a comfortable team in which nobody is
challenged. It is one in which a person can surface bad news, question a
forecast, report harassment, admit an error, or oppose the founder without
retaliation, while claims are still examined and decisions are still made.

Leaders can test the culture with four questions:

\begin{enumerate}[itemsep=0.14em]
\item What truth becomes expensive for an employee to say here?
\item Which producer, founder, customer, or investor is treated as exempt from
  the rules?
\item Who bears the hidden cost when a target is met?
\item When someone last stopped unsafe or misleading work, what happened?
\end{enumerate}

The answers reveal the operating constitution more reliably than the values
page.

Trust and controls perform different work. Trust lowers the cost of
cooperation and makes unwelcome truth speakable. Stable rules, records, and
independent review keep a favored producer or founder from purchasing
immunity. Personal trust alone invites favoritism and silence; surveillance
alone teaches people to game the visible measure and conceal ambiguity. A
credible company states the rule and its reason, examines proportionately, and
leaves a legitimate path for challenging both the claim and the rule.

\section{Do not make another person carry invisible risk}

People extend a company's capacity only if the company accepts duties in
return. Those duties include lawful and understandable pay, safe tools and
staffing, accurate expectations, reasonable workload, protection from
discrimination and retaliation, credit for contribution, access to relevant
information, and a way to raise concerns. Where people create durable value,
shared upside may also be appropriate. Keller's injunction to help talented
employees get rich is not a compensation formula, but it correctly rejects the
idea that all gains naturally belong to the person whose name is on the shares.

One real-estate operator recounts an incident in which his pregnant wife was
chased with a stick and a dog while visiting a managed property. He presents it
as one of the strange events that had begun to feel normal. The useful lesson
is not the spectacle. Repeated exposure can normalize danger until the absence
of catastrophe is mistaken for a safe process.

Field work needs a hazard assessment before dispatch: known occupants and
conditions, check-in and exit procedures, whether two people or trained
security are needed, communication, de-escalation limits, emergency support,
and an unconditional right to leave. No revenue target makes an avoidable
threat part of the job. After an incident, protect the person, preserve facts,
report where required, correct the system, and do not reward the manager for
having completed the visit anyway.

Fairness also improves the quality of decisions. A worker who cannot afford to
lose a shift may hide an injury. A manager whose bonus depends only on volume
may suppress a defect. A partner excluded from accounts cannot challenge a
cash call. Compensation, measurement, information, and voice are therefore
governance mechanisms, not human-resources decoration.

\section{A leader needs somewhere to put the truth}

The Austin real-estate founder describes leading while deals collapse,
complaints arrive, and lawsuits demand attention. His instinct is not to let
his own struggle show. A leader does owe the organization steadiness; anxiety
should not be broadcast merely to transfer it. But performed invulnerability
creates a second danger. If nobody can see pressure, nobody can help assess it,
and the leader begins to confuse concealment with protection.

Jim Keyes, reflecting on corporate crises, joins cash, communication, and
character. Communication matters because a frightened team cannot execute a
direction it does not understand. The answer is calibrated truth: disclose the
situation, what is known and unknown, what changes now, what remains stable,
who decides next, and when the team will hear more. Do not promise certainty
that does not exist.

The leader also needs a smaller place for the unprocessed truth: a board member,
co-founder, qualified adviser, peer, clinician, spouse, or other trusted person
whose role and confidentiality fit the problem. That support should challenge
judgment rather than merely applaud it. Legal, financial, safety, and mental-
health problems require relevant competence. The organization should never
depend on one person's ability to absorb fear without consequence.

\section{Update Record 3 with the people-and-governance charter}

Choose the smallest real operating unit that matters now: a partnership, a
leadership team, a field crew, or the people responsible for one customer
promise. Write a charter short enough to use and specific enough to settle a
hard day. Add six connected fields to Record 3:

\begin{itemize}[leftmargin=0pt,label={},labelsep=0pt,itemsep=0.18em,
  parsep=0pt,topsep=0.35em]
\item \textbf{Promise and exposure.} State the purpose, customer promise, and
  boundaries the unit will not cross for growth. Name the customers, workers,
  owners, creditors, suppliers, communities, and other people who carry a
  material consequence.
\item \textbf{Responsibilities.} Give each person the eight-field bounded
  responsibility used in this chapter. Record the continuing contribution,
  direct customer signal, and reversible evidence required before equity or a
  wider consequence.
\item \textbf{Economics.} Separate labor, management, capital, relationships,
  intellectual property, and personal guarantees. Record compensation,
  equity, vesting or repurchase, distributions, expenses, benefits, credit,
  and future capital calls.
\item \textbf{Decisions and truth.} Assign ordinary domains, consultation,
  final closure, spending limits, information rights, and the treatment of
  dissent. Define the trigger, scope, communication, evidence, review, and
  expiry of emergency authority.
\item \textbf{Protection and support.} Preserve workload, staffing, safety,
  privacy, anti-harassment, and non-retaliation duties through a usable
  reporting and response path. Name competent support for the work and a
  confidential place where leaders can test unprocessed pressure.
\item \textbf{Repair and continuity.} Define feedback, coaching, correction,
  immediate and ordinary separation, conflict, independent review, deadlock,
  incapacity, misconduct, departure, transfer, valuation, dissolution, and
  sale. Name the remedy after harm and a quarterly review of whose
  contribution, risk, or voice remains invisible.
\end{itemize}

Have the people governed by the charter examine it. Obtain qualified legal,
tax, employment, insurance, and safety advice where the stakes require it.
Then test the document against a late customer payment, a failed launch, a
harassment report, a partner who wants to leave, a founder who becomes ill, and
a star producer who lies. If the answer is still ``the founder will decide,''
the architecture is unfinished.

A company becomes more than its founder when other people can contribute
judgment without surrendering dignity and can carry responsibility without
being denied authority. That achievement still leaves a harder ownership
question. The people may operate the machine, but whoever controls its critical
rights, data, capital, and routes to the customer can still control its
economics. We turn next to those points of dependence.
